% !TEX root = ../Main.tex
\chapter{三角换元}

三角换元的根基是那条最基本的恒等式
\[
\sin^2\theta+\cos^2\theta=1.
\]
凡是题目里出现"平方和"或"圆"的结构，就可以用一个角 $\theta$ 把两个变量串起来，把约束自动满足掉。

\section{$\sin\theta$ 和 $\cos\theta$ 的再定义}

在\textbf{单位圆} $x^2+y^2=1$ 上取一点 $(x,y)$，设它与原点连线和 $x$ 轴正方向的夹角为 $\theta$，则规定
\[
\cos\theta=x,\qquad \sin\theta=y.
\]

\begin{center}
\begin{tikzpicture}[>=Stealth,scale=1.4]
  \draw[->] (-1.4,0) -- (1.5,0) node[right]{$x$};
  \draw[->] (0,-1.4) -- (0,1.5) node[above]{$y$};
  \draw[thick,blue!60!black] (0,0) circle (1);
  \coordinate (P) at (40:1);
  \draw[thick] (0,0) -- (P);
  \draw (0.32,0) arc (0:40:0.32);
  \node at (0.5,0.18) {\footnotesize$\theta$};
  \fill (P) circle (0.6pt) node[above right]{\footnotesize$(x,y)$};
  \draw[dashed] (P) -- (40:1 |- 0,0);
\end{tikzpicture}
\end{center}

让 $\theta$ 沿圆周扫一圈，点 $(x,y)$ 始终在圆上，于是 $x^2+y^2=1$ 永远成立，也就\textbf{反过来得到}
\[
\cos^2\theta+\sin^2\theta=1.
\]
这就是用几何（单位圆）重新认识 $\sin,\cos$。

\section{平方和为 1 的代换}

只要看到形如
\[
(x-1)^2+(y-2)^2=1
\]
的"两个平方相加等于 $1$"，就把它和 $\cos^2\theta+\sin^2\theta=1$ 对照：
\[
\underbrace{(x-1)^2}_{\cos^2\theta}+\underbrace{(y-2)^2}_{\sin^2\theta}=1
\quad\Longrightarrow\quad
\bc
x-1=\cos\theta,\\
y-2=\sin\theta,
\ec
\]
即
\[
\bc
x=1+\cos\theta,\\
y=2+\sin\theta.
\ec
\]
两个变量 $x,y$ 被压缩成一个角 $\theta$，约束自动满足。

\section{圆的三角代换}

把上面的想法推广到一般的圆。单位圆 $x=\cos\theta,\ y=\sin\theta$ 经过两步变换就能覆盖任意圆：

\begin{itemize}
    \item \textbf{放缩}：半径从 $1$ 放大到 $R$，得 $(R\cos\theta,\ R\sin\theta)$；
    \item \textbf{平移}：圆心从 $(0,0)$ 移到 $(a,b)$。
\end{itemize}
于是圆 $(x-a)^2+(y-b)^2=R^2$ 上的点可写成
\[
\bc
x=a+R\cos\theta,\\
y=b+R\sin\theta.
\ec
\]

\section{直线的三角代换}

直线也能用同样的"方向角"思想参数化。设直线过定点 $P(a,b)$，倾斜角为 $\theta$，$Q(x,y)$ 是线上动点、到 $P$ 的有向距离为 $t$。

\begin{center}
\begin{tikzpicture}[>=Stealth,scale=1.2]
  \draw[->] (-0.5,0) -- (3,0) node[right]{$x$};
  \draw[->] (0,-0.4) -- (0,2.4) node[above]{$y$};
  \coordinate (P) at (0.7,0.6);
  \coordinate (Q) at (2.3,1.8);
  \draw[thick,blue!60!black] (0.2,0.225) -- (2.6,2.025);
  \draw[dashed] (Q) -- (Q |- P) node[midway,right]{\footnotesize$t\sin\theta$};
  \draw[dashed] (P) -- (Q |- P) node[midway,below]{\footnotesize$t\cos\theta$};
  \draw (P) ++(0.45,0) arc (0:37:0.45);
  \node at ($(P)+(0.65,0.18)$) {\footnotesize$\theta$};
  \fill (P) circle (0.8pt) node[below left]{\footnotesize$P(a,b)$};
  \fill (Q) circle (0.8pt) node[above right]{\footnotesize$Q(x,y)$};
\end{tikzpicture}
\end{center}

把位移 $t$ 沿 $x,y$ 两个方向分解，得到直线的参数式
\[
\bc
x=a+t\cos\theta,\\
y=b+t\sin\theta.
\ec
\]
这里 $\theta$ 是固定的（直线方向），变的是 $t$；它和圆的代换形式上一致，区别只是圆里 $R$ 固定、$\theta$ 变，直线里 $\theta$ 固定、$t$ 变。
